% Original contribution drafted 25 November 2026.
% Intended as an input section in the main report. Requires amsmath, amsthm.
% No novelty priority or full Lean verification is claimed.
\section{An exact torsion specialization with a uniform order bound}
\label{sec:exact-torsion}

The preceding quantitative report retained a positive valuation tolerance.
The argument below reaches tolerance zero for one polynomial, and hence
for the aggregate valuation of a finite packet.  Its additional input is
discrete: a nonempty multiset of roots on a fixed cyclotomic shell has a
definite minimum mass.  The old-shell saving is retained exactly instead
of replacing it everywhere by a fixed fraction of accumulated mass.

Fix a prime $p$ and the valuation $v$ on $\mathbf C_p$, normalized by
$v(p)=1$, with $v(0)=+\infty$.  Let $\mathcal O$ be its valuation ring.
For $h\geq1$ write
\[
 e_h=(p-1)p^{h-1},\qquad r_h=e_h^{-1},\qquad
 U_h=\{\zeta:\operatorname{ord}(\zeta)=p^h\},\qquad
 q=p^{-1},\quad c=1-q.
\]
For a nonzero integral polynomial $g(T)=\sum_j a_jT^j$ define
\[
 G_\rho(g)=\min_j\{v(a_j)+j\rho\}\qquad(\rho\geq0).
\]
Thus $G_0(g)\geq0$, and $v(g(x))\geq G_\rho(g)$ whenever
$v(x)\geq\rho$.  All logarithmic valuations below are additive.

\begin{theorem}[Uniform exact specialization]\label{thm:exact-selector}
Let $\rho>0$, $B\geq0$, and $h_0\geq1$ satisfy $r_{h_0}\leq\rho$.
Put $a=r_{h_0}$ and
\[
 H_0=h_0+\left\lceil p^{h_0}B(h_0+B+2)\right\rceil.
\]
For every nonzero $g\in\mathcal O[T]$ with $A=G_\rho(g)\leq B$,
there exist $h_0\leq h\leq H_0$ and $\zeta\in U_h$ such that
\[
                 v(g(\zeta-1))\leq G_\rho(g).
\]
If $B<a$, every $\zeta\in U_{h_0}$ already has this property.
The order bound is independent of degree, coefficients, and their field
of definition.  It is an upper bound on an available order exponent;
it does not specify that $h=H_0$.
\end{theorem}

Here $(1-1/p)r_{h_0}=p^{-h_0}$.  For fixed $p,h_0$, this bound is
quadratic in $B$ and contains no
tolerance parameter.  The estimate is deliberately conservative and is
not asserted to be optimal.

\subsection{The shell calculation}

Factor $g=b\prod_\alpha(T-\alpha)$, counting multiplicities.  Gauss
multiplicativity gives
\[
 G_u(g)=v(b)+\sum_\alpha\min(u,v(\alpha)).
\]
This formula includes roots of negative valuation and roots at zero.
Let $R_h$ be the multiset of roots with valuation exactly $r_h$, and set
\[
 w_h=r_h\#R_h,\qquad
 W_h=\sum_{k=h_0}^h w_k,\qquad
 z_h=\sum_{k=h_0}^h q^{h-k}w_k,
 \qquad W_{h_0-1}=z_{h_0-1}=0.
\]
For any finite endpoint $H$,
\begin{equation}\label{eq:root-budget-exact}
 W_H\leq G_\rho(g)-G_0(g)\leq A\leq B.
\end{equation}
Indeed, every root counted on these shells contributes $r_k$ to the
difference of the two Gauss valuations.  At shell $h$, put
$s_h=A-G_{r_h}(g)\geq0$.  The contribution of earlier shells yields
\begin{equation}\label{eq:exact-old-shell-slack}
 s_h\geq\sum_{k=h_0}^{h-1}(r_k-r_h)\#R_k
      = W_{h-1}-qz_{h-1}.
\end{equation}
The factor $q^{h-k}$ is retained here.

For $\alpha\in R_h$ and $\zeta\in U_h$, let
$E_\alpha(\zeta)=v(\zeta-1-\alpha)-r_h\geq0$.
Roots outside $R_h$ have no cancellation with $\zeta-1$, so
\begin{equation}\label{eq:exact-value-shell}
 v(g(\zeta-1))=G_{r_h}(g)+\sum_{\alpha\in R_h}E_\alpha(\zeta).
\end{equation}

\begin{lemma}[Capped average]\label{lem:capped-average-exact}
For $C>0$ and $\alpha\in R_h$,
\[
 \frac1{e_h}\sum_{\zeta\in U_h}\min(E_\alpha(\zeta),C)
 \leq r_h(h+C).
\]
\end{lemma}
\begin{proof}
Choose a primitive root $\zeta_*$ maximizing
$v(\zeta-(1+\alpha))$.  For $\zeta\ne\zeta_*$ the ultrametric
inequality gives $v(\zeta-(1+\alpha))\leq v(\zeta-\zeta_*)$.
The term at $\zeta_*$ is at most $C$.  The sum of the valuations of
the remaining differences is
\[
 v(\Phi'_{p^h}(\zeta_*))=h-\frac1{p-1}\leq h,
\]
obtained by differentiating
$\Phi_{p^h}(X)=(X^{p^h}-1)/(X^{p^{h-1}}-1)$ at $\zeta_*$.
Discarding the subtractive $r_h$ terms and dividing by $e_h$ proves
the stated bound.  This also covers $p=2,h=1$, when the set has one
element and the derivative valuation is zero.
\end{proof}

More precisely, the right side can be replaced by
$r_h(C+h-1-1/(p-1)+r_h)$; that improvement is not needed.

Suppose every primitive root on shell $h$ violates
$v(g(\zeta-1))\leq A+\delta$, where $\delta\geq0$.
Choose $C>B+\delta$.  In view of \eqref{eq:exact-value-shell}, the
sum of the excesses is greater than $s_h+\delta$.  Capping each
excess at $C$ preserves this strict inequality, because
$s_h+\delta\leq B+\delta<C$.  Averaging and applying the lemma gives
\begin{equation}\label{eq:failure-recurrence-input}
 w_h>\frac{Q_h}{h+C},\qquad
 Q_h=\delta+W_{h-1}-qz_{h-1}.
\end{equation}
The states satisfy
\begin{equation}\label{eq:state-update}
 z_h=qz_{h-1}+w_h,\qquad Q_{h+1}=Q_h+c z_h.
\end{equation}
In particular, all entries of the state-update matrix obtained by
putting $w_h=Q_h/(h+C)$ are nonnegative.  Replacing the strict lower
bound by equality therefore gives a componentwise lower comparison
sequence for the actual states and weights.

\subsection{The discrete seed and proof of the exact theorem}

Take $\delta=0$ and $C=B+1$.  If all shells from $h_0$ through $H$
are bad, then $R_{h_0}\ne\varnothing$: otherwise
\eqref{eq:exact-value-shell} would give
$v(g(\zeta-1))=G_{r_{h_0}}(g)\leq A$ on the first shell.  Thus
\[
                  w_{h_0}\geq a=r_{h_0}.
\]
This proves the assertion for $B<a$ immediately.

For subsequent shells compare with the equality recurrence starting
from $z_{h_0}=a$ and $Q_{h_0+1}=ca$.  Write $u_h=Q_h/t_h$ for its
weights, where $t_h=h+C$, and put
$d_h=u_h-cz_{h-1}$.  Direct substitution into
\eqref{eq:state-update} gives
\begin{equation}\label{eq:defect-recurrence}
 u_{h+1}-u_h=-\frac{q\,d_h}{t_h+1},\qquad
 d_{h+1}=q\frac{t_h}{t_h+1}d_h.
\end{equation}
At $h=h_0+1$,
\[
 u_{h_0+1}=\frac{ca}{h_0+B+2},\qquad
 d_{h_0+1}=ca\left(\frac1{h_0+B+2}-1\right)<0.
\]
Hence every later $d_h$ is negative and the equality weights $u_h$
increase.  Actual weights dominate this comparison sequence, so
\[
 W_H\geq a+(H-h_0)\frac{ca}{h_0+B+2}.
\]
For $H=H_0$, using $ca=p^{-h_0}$, the right side is at least $a+B>B$, contradicting
\eqref{eq:root-budget-exact}.  The case $B=0$ is already covered by
$B<a$.  This proves Theorem~\ref{thm:exact-selector}.

\subsection{A sharper positive-tolerance estimate}

The same recurrence also improves the previous polynomial dependence
on inverse tolerance.

\begin{theorem}[Linear inverse-tolerance bound]\label{thm:linear-tolerance}
In the notation of Theorem~\ref{thm:exact-selector}, fix $\delta>0$,
and define
\[
 t_0=h_0+B+\delta+1,\qquad
 \kappa=1-\frac1{(p-1)(t_0+1)}>\frac12,
 \qquad
 H_+=h_0+\left\lceil\frac{B t_0}{\kappa\delta}\right\rceil.
\]
Some $h_0\leq h\leq H_+$ and $\zeta\in U_h$ satisfy
$v(g(\zeta-1))\leq G_\rho(g)+\delta$.
\end{theorem}
\begin{proof}
Use $C=B+\delta+1$ in \eqref{eq:failure-recurrence-input} and
compare with the equality recurrence starting at
$Q_{h_0}=\delta$, $z_{h_0-1}=0$.  Now
$u_{h_0}=d_{h_0}=\delta/t_0$, and
\eqref{eq:defect-recurrence} gives
\[
 d_{h_0+j}=\frac{\delta q^j}{t_0+j},\qquad
 u_{h_0+j}=\frac\delta{t_0}
  -\delta\sum_{k=0}^{j-1}
       \frac{q^{k+1}}{(t_0+k)(t_0+k+1)}
 \geq\frac{\kappa\delta}{t_0}.
\]
The last inequality bounds each denominator below by $t_0(t_0+1)$
and sums the geometric series.  If every shell through $H_+$ were
bad, \eqref{eq:root-budget-exact} would imply
$B\geq W_{H_+}\geq(H_+-h_0+1)\kappa\delta/t_0>B$.
\end{proof}

For fixed $p,h_0,B$ and $0<\delta\leq1$, $H_+=O(1/\delta)$.
For very small $\delta$, the zero-tolerance bound $H_0$ is stronger.
Neither bound is a polynomial bound for the root order $p^H$ itself.

\subsection{Multivariate and bounded analytic versions}

For a polynomial or a bounded integral formal power series
$F(\mathbf T)=\sum_\nu a_\nu\mathbf T^\nu$ in $s<\infty$ variables,
write
\[
 G_\rho(F)=\min_\nu\{v(a_\nu)+\rho|\nu|\}.
\]
For nonzero series and $\rho>0$ the minimum exists: the integral
coefficients and $\rho|\nu|\to\infty$ leave only finitely many
possible minimizing terms.  The series converges wherever every
coordinate has positive valuation.

\begin{corollary}[Exact analytic specialization]\label{cor:exact-analytic}
The bound $H_0$ in Theorem~\ref{thm:exact-selector} works for every
finite number $s$ of variables and every nonzero
$F\in\mathcal O[[T_1,\ldots,T_s]]$ with $G_\rho(F)\leq B$.
There are $\zeta_j\in U_{h_j}$, $h_0\leq h_j\leq H_0$, such that
\[
 v(F(\zeta_1-1,\ldots,\zeta_s-1))\leq G_\rho(F).
\]
In particular, one may replace $G_\rho(F)$ on the right by a finite
witness valuation $v(F(\mathbf x))$ with $v(x_j)\geq\rho$.
The bound is independent of $s$; the available finite grid has size
that grows with $s$.
\end{corollary}
\begin{proof}
First take $F$ polynomial.  Write
$F=\sum_\beta f_\beta(T_1)\mathbf T_{\rm rest}^\beta$ and choose
$\beta$ attaining
$G_\rho(F)=G_\rho(f_\beta)+\rho|\beta|$.
Since $0\leq G_\rho(f_\beta)\leq B$, the scalar theorem chooses
$\zeta_1$ with
$v(f_\beta(\zeta_1-1))\leq G_\rho(f_\beta)$.
The remaining polynomial after substitution has integral coefficients,
is nonzero, and has Gauss valuation no larger than $G_\rho(F)$.
Repeat.  This loses no tolerance, so the same $H_0$ serves each step.

For a series, set $r_*=r_{H_0}>0$ and choose $D$ with
$(D+1)\min(\rho,r_*)>B$.  The truncation $P$ to total degree at most
$D$ has $G_\rho(P)=G_\rho(F)$.  At every point of the finite torsion
grid in the corollary, the tail $F-P$ has valuation greater than $B$.
The polynomial theorem gives $v(P)\leq B$ there, so the ultrametric
inequality forces $v(F)=v(P)$.  Integral coefficients in $\mathbf C_p$
are permitted; no algebraicity assumption is needed in the polynomial
root-factorization argument.
\end{proof}

\begin{corollary}[Exact aggregate budget for a packet]\label{cor:aggregate}
Let $F_1,\ldots,F_n\in\mathcal O[[T_1,\ldots,T_s]]$ be nonzero,
and assume $\sum_iG_\rho(F_i)\leq B$.
There is a common torsion tuple with exponents in $[h_0,H_0]$ such
that
\[
        \sum_{i=1}^n v(F_i(\boldsymbol\zeta-\mathbf1))
        \leq\sum_{i=1}^nG_\rho(F_i)\leq B.
\]
The same conclusion follows from separate witness points
$\mathbf x_i$ with coordinates of valuation at least $\rho$ and
$\sum_i v(F_i(\mathbf x_i))\leq B$.
\end{corollary}
\begin{proof}
Apply Corollary~\ref{cor:exact-analytic} to $\prod_iF_i$ and use
Gauss multiplicativity.  The witness version follows from
$G_\rho(F_i)\leq v(F_i(\mathbf x_i))$ separately for every $i$.
All evaluated factors are integral, so their valuations are nonnegative;
the finite bound implies that none vanishes.
\end{proof}

This is an aggregate statement.  It does not assert
$v(F_i(\boldsymbol\zeta-\mathbf1))\leq G_\rho(F_i)$ for each $i$.
For example, negative deviations of other factors from their original
Gauss values can compensate for a positive deviation in one factor.
This distinction matters for arithmetic applications that impose
separate bounds on individual Selmer groups.

\begin{proposition}[No cardinality-free exact individual bounds]
\label{prop:individual-obstruction}
Fix $p=2$, $h_0\geq1$ and $\rho\geq r_{h_0}$.  For every $L\geq h_0$
there is a finite packet with total Gauss budget strictly less than
$2r_{h_0}$ for which no primitive root on any shell
$h_0,\ldots,L$ satisfies all the exact individual Gauss bounds.
\end{proposition}
\begin{proof}
Choose one primitive $2^h$-th root $\xi_h$ on each shell and set
$g_h(T)=T-(\xi_h-1)$.  Then
$G_\rho(g_h)=r_h=v(g_h(0))$, and
$\sum_{h=h_0}^Lr_h<2r_{h_0}$.
For any primitive $2^h$-th root $\zeta$, either $\zeta=\xi_h$, or
$\zeta/\xi_h$ has order at most $2^{h-1}$ because both roots have
odd exponents relative to a primitive $2^h$-th root.  Thus
$v(g_h(\zeta-1))=v(\zeta-\xi_h)>r_h$; at $h=1$ the value is zero
and its valuation is infinite.  Constraint $h$ fails throughout shell
$h$, proving the assertion.
\end{proof}

The same construction also explains why the radius in the exact
theorem must remain the fixed comparison radius $\rho$.

\begin{proposition}[The moving Gauss target has no such uniform bound]
\label{prop:moving-gauss-obstruction}
For $p=2$ and every $L\geq1$, an integral polynomial $g$ with
$G_1(g)<2$ can satisfy
$v(g(\zeta-1))>G_{r_h}(g)$ for every $1\leq h\leq L$ and every
$\zeta\in U_h$.
\end{proposition}
\begin{proof}
Take $g(T)=\prod_{h=1}^L(T-(\xi_h-1))$ as above.
Its budget is $G_1(g)=\sum_{h=1}^Lr_h=2-2^{1-L}<2$.
On shell $h$, roots on other shells contribute exactly their
Gauss terms, whereas the factor indexed by $h$ contributes strictly
more than $r_h$.  Formula~\eqref{eq:exact-value-shell} gives the claim.
\end{proof}

The upstream positive-tolerance lemma compares against the moving
$G_{r_h}$ and remains a genuinely stronger comparison in that respect.

For an arbitrary family with finite total Gauss budget, a single point
also satisfies the aggregate bound.  Apply the finite-packet assertion
to each finite subfamily.  For the fixed finite torsion grid, the sets
where a finite subfamily has total valuation at most the global budget
have the finite intersection property, since valuations are nonnegative.
Their intersection is nonempty.  This is a finite-grid compactness
argument, not a convergence claim for an infinite product.

\paragraph{First powers.}
Let $m\geq1$ and suppose the packet has separate witnesses with total
valuation at most $B$.  For $1\leq j\leq m$, form
$F_i((1+T_1)^j-1,\ldots,(1+T_s)^j-1)$.
Each transformed series has Gauss valuation at radius $\rho/m$ at most
the original witness value.  Indeed, for each coordinate one can choose
$y$ with $(1+y)^j=1+x$ and $v(y)\geq v(x)/j\geq\rho/m$.
For $x\ne0$, this follows by factoring the monic polynomial
$(1+Y)^j-1-x$: its roots are integral and the sum of their valuations
is $v(x)$; for $x=0$ choose $y=0$.
Applying the aggregate corollary at radius $\rho/m$ and budget $mB$
therefore gives one tuple with
\[
 \sum_{i=1}^n\sum_{j=1}^m
 v(F_i(\zeta_1^j-1,\ldots,\zeta_s^j-1))\leq mB.
\]
All values are nonzero.  This is neither a prescribed-order theorem nor
a density theorem for arithmetic characters of bounded conductor.

\subsection{Limitations and comparison with the previous report}

The earlier report proves a positive fraction of good primitive roots
and gives separate per-polynomial tolerances.  Those conclusions are
not superseded by the exact aggregate theorem.  The improvement here is
the tolerance-zero endpoint and its explicit uniform bound for a single
polynomial or product, together with the linear inverse-tolerance scalar
bound.  It supplies a local specialization ingredient.  A transfer to a
particular $L$-function still requires integral interpolation data and
nonzero witnesses; a transfer to Selmer groups requires the repository's
representation hypotheses and normalized-length bookkeeping.  No named
major conjecture or worldwide novelty claim follows from the present
argument alone.
